\emph{(i).} By Theorem \claimref{hawking.greybody-inheritance-under-interface-premise} the flux in each channel is $\frac{1}{2\pi}\Gamma_{s\ell m}(E)\,n_{\pm}(E)$ at $\TH = 1/(8\pi M)$. For the Schwarzschild exterior the transmission coefficients depend on $E$ only through $ME$ (the barriers scale with $M$), and $E/\TH = 8\pi ME$, so each channel's flux per unit energy is $\phi(ME)$ for a function $\phi$ fixed by the channel. The power is then $P(M) = \int E\,\phi(ME)\,dE = M^{-2}\int x\,\phi(x)\,dx$ (substituting $x = ME$), and summing over the channels of every populated tower, which by Proposition \claimref{dark-matter.parallel-boundary-mode-spaces-one-wall}(i) all radiate from the one exterior in proportion to their light degrees of freedom once the temperature exceeds their gaps, gives $dM/dt = -\Psi_{\mathrm{tot}}/M^{2}$ with constant $\Psi_{\mathrm{tot}}$ in the stated regime. Integrating to the end of the evaporation, $M(\tau)^{3} = 3\Psi_{\mathrm{tot}}\tau$, and at fixed $M$ this tower's power $\Psi_{\mathrm{vis}}/M^{2}$ is the same as in the standard case. The check verifies the $M^{-2}$ law exactly for polynomial flux shapes of compact support in $ME$.

\emph{(ii).} Let $M_{\Psi}(t)$ be the mass at remaining lifetime $t$ when the total coefficient is $\Psi$, so $M_{\Psi_{\mathrm{tot}}}(t)^{3} = 3\Psi_{\mathrm{tot}}t = 3\Psi_{0}(Dt) = M_{\Psi_{0}}(Dt)^{3}$. The visible emission per unit time at mass $M$ is a function $h(E, M)$ of the exterior and of this tower's content only: Proposition \claimref{dark-matter.parallel-boundary-mode-spaces-one-wall}(i) states that barriers, transmission coefficients and occupancies do not record which tower the mass came from, and the same holds for the visible secondaries, which are this tower's physics at energy $E$. Hence
\[
N(E;\tau) = \int_{0}^{\tau} h\bigl(E, M_{\Psi_{\mathrm{tot}}}(t)\bigr)\,dt = \int_{0}^{\tau} h\bigl(E, M_{\Psi_{0}}(Dt)\bigr)\,dt = \frac{1}{D}\int_{0}^{D\tau} h\bigl(E, M_{\Psi_{0}}(s)\bigr)\,ds = \frac{1}{D}N_{0}(E;D\tau),
\]
with $s = Dt$. The check verifies the identity exactly for emission functions polynomial in $M^{3}$, both with the baseline $\Psi_{0} = \Psi_{\mathrm{vis}}$ and with a graviton term in it, and confirms that the plain ratio $N_{0}(E;\tau)/D$ differs from it.

\emph{(iii).} The visible energy emitted in the last $\tau$ is the visible share of the remaining mass, $(\Psi_{\mathrm{vis}}/\Psi_{\mathrm{tot}})M_{\Psi_{\mathrm{tot}}}(\tau) = D^{-2/3}(\Psi_{\mathrm{vis}}/\Psi_{0})M_{\Psi_{0}}(\tau)$, $D^{-2/3}$ times the standard value. The temperature is $1/(8\pi M_{\Psi_{\mathrm{tot}}}(\tau)) = D^{-1/3}/(8\pi M_{\Psi_{0}}(\tau))$. For the primary visible flux, $h(E, M) = \phi_{\mathrm{vis}}(ME)$ by (i); with $dt = M^{2}dM/\Psi_{\mathrm{tot}}$,
\[
N(E;\tau) = \frac{1}{\Psi_{\mathrm{tot}}}\int_{0}^{M(\tau)}\phi_{\mathrm{vis}}(ME)\,M^{2}\,dM = \frac{1}{\Psi_{\mathrm{tot}}E^{3}}\int_{0}^{EM(\tau)}\phi_{\mathrm{vis}}(x)\,x^{2}\,dx,
\]
which for $EM(\tau)$ large tends to $C E^{-3}/\Psi_{\mathrm{tot}} = C E^{-3}/(D\Psi_{0})$ with $C = \int_{0}^{\infty}\phi_{\mathrm{vis}}(x)x^{2}\,dx$, finite since the occupancies fall exponentially in $x$, and independent of $\tau$: the tail is $1/D$ times the standard tail. Secondaries obey the exact rescaling (ii) but not this tail law, which needs $h$ to depend on $E$ and $M$ only through $ME$. The check verifies the cube relations exactly and the $\tau$-independence of the tail for a flux shape of compact support. At fixed $M$ the visible power is unchanged and the visible fraction of the total is $1/D$ of its standard value, so the full ratio applies to that fraction and to the primary tail, and not to quantities measured at fixed time before the end.

\emph{(iv).} By Theorem \claimref{dark-matter.mirror-tower-copy-one}(i) and (iii) the mirror's light content, and hence its emission coefficient, equals this tower's, so $\Psi_{\mathrm{tot}} = 2\Psi_{\mathrm{vis}} + \Psi_{\mathrm{odd}} + \Psi_{\mathrm{grav}}$ and $D = (2 + g + o)/(1 + g) = 2 + (o - g)/(1 + g)$, which is at least $2$ exactly when $o \ge g$; with $o = 0$ it is $2 - g/(1 + g)$, just below $2$ since $g$ is small for the standard-model content (of order a percent or less, from the literature, cited from memory). The odd towers' content (Propositions \claimref{dark-tower-matter.three-dimensional-content} and \claimref{dark-tower-matter.five-dimensional-content}) enters only through $o$, whose value needs the emission of lower- and higher-dimensional fields from the one exterior, which is not derived here.

\emph{Numerical reading (not part of the check).} Burst searches compare the photon fluence of the final $\tau$, for $\tau$ between a fraction of a second and a few hundred seconds, with the standard expectation; their energy band lies near the temperature of the last seconds (a few TeV and above for a hole of the standard-model content, from the literature, cited from memory), where neither the tail law nor the energy law is exact. At $D = 2$ the fluence there is dimmed by a factor between $2^{2/3} \approx 1.59$ and $2$ (computed), and for Euclidean, background-free counting the sensitive volume, which scales as the fluence to the power $3/2$, falls by a factor between $2$ and $2^{3/2} \approx 2.83$; so a rate limit derived for undimmed bursts weakens by that factor.
